\chapter{Conclusões Parciais}
\label{cap:05}

Este trabalho tem como objetivo desenvolver um protótipo de ferramenta computacional para a análise comparativa de proteínas associadas a vias metabólicas do sabor e do aroma do café. Até a qualificação, foram concluídas as cinco primeiras etapas da metodologia, que compõem o núcleo bioinformático, e parte da validação. Os cinco primeiros objetivos específicos foram, portanto, atendidos, o sétimo foi atendido em parte, e o sexto, referente à API e à interface web, conta apenas com um protótipo visual.

O núcleo localizou e comparou os genes \textit{DXMT}, \textit{HCT} e \textit{SPS} em quatro proteomas do gênero \textit{Coffea}. O HCT e a SPS foram encontrados nos quatro, com identidades de 100\% entre as variedades Geisha e Red Bourbon de \textit{C. arabica}, de 99,52\% a 100\% entre estas e \textit{C. canephora}, e de 92,56\% a 97,70\% nas comparações com \textit{C. mauritiana}. Esses valores acompanham o grau de parentesco entre os organismos comparados e são coerentes com a conservação entre variedades de \textit{C. arabica} relatada na literatura, o que indica que o fluxo implementado produz resultados plausíveis.

Três conclusões decorrem do que foi realizado. A primeira é que o melhor resultado de uma busca por similaridade não basta para identificar o gene correspondente em famílias com duplicações recentes. No caso do DXMT, a busca sem confirmação recíproca selecionou, em dois proteomas, um parálogo do agrupamento de N-metiltransferases, e os valores de identidade resultantes estavam incorretos; a confirmação por melhores correspondências recíprocas eliminou o erro. A segunda é que a implementação própria do algoritmo de Gotoh é confiável para o uso pretendido, pois suas pontuações coincidiram com as de uma implementação de referência em todos os casos testados. A terceira é que a qualidade dos dados de entrada condiciona o alcance da comparação: dois dos seis proteomas obtidos foram descartados por apresentarem menor integridade, e as comparações do DXMT falharam em dois dos quatro restantes, na Geisha por falta de um modelo gênico na anotação e em \textit{C. mauritiana} por não haver cópia íntegra do gene.

Os resultados têm limitações que devem ser consideradas. O conjunto analisado é pequeno, com três genes e quatro proteomas, e não sustenta generalizações sobre as vias metabólicas. Os proteomas de \textit{C. arabica} utilizados contêm apenas o subgenoma de origem \textit{canephora}, o que restringe a comparação a um dos dois subgenomas da espécie. Além disso, os valores obtidos descrevem a conservação de sequências e não permitem inferir diferenças de sabor ou de aroma entre as variedades.

As etapas seguintes, previstas no cronograma (Quadro~\ref{qua:cronograma}), são a resolução das pendências do núcleo, com o registro das versões dos dados; o desenvolvimento da API e da interface web, a partir do protótipo visual apresentado; e a conclusão da validação, com a conferência dos genes localizados em relação aos identificadores publicados.
